% Save the current style
\let\origthefootnote\thefootnote
\renewcommand{\thefootnote}{\fnsymbol{footnote}}

\epigraph{
    This \textup{[\(e^z\)]} is the most important function in mathematics.
}{
    % Real and Complex Analysis\\
    \textsc{---Walter Rudin}\footnotemark
}

\footnotetext{\fullcite{Rudin_1987}}

% Restore original footnote formatting
\let\thefootnote\origthefootnote

\section{Uvod} \label{sec:intro}

Naj bo \(x_0\) poljubno realno število. Oglejmo si, kaj se dogaja z zaporedjem iteracij števila \(x_0\) glede na funkcijo \(e^{x}\):
\[x_0 \mapsto e^{x_0} \mapsto e^{e^{x_0}} \mapsto e^{e^{e^{x_0}}} \mapsto \cdots\]
Hitro se prepričamo, da je vsako tako zaporedje divergentno. Res, zaporedje zgornjih števil je ne glede na začetni pogoj \(x_0\) strogo naraščajoče in neomejeno, njegova dinamika pa razmeroma nezanimiva.

Povsem drugačno obnašanje eksponentna funkcija ponudi, če jo s predpisom \(e^{x + iy} = e^x (\cos y + i \sin y)\) razširimo na kompleksno ravnino. Izkaže se, da na vsaki odprti množici obstajajo števila, katerih orbite se med seboj kvalitativno razlikujejo. Analiza njihovega obnašanja bo osrednji problem tega diplomskega seminarja. Opazovali bomo torej rekurzivno zaporedje, dano s predpisom \(z_{n + 1} = e^{z_n}\). Ogledali si ga bomo za različne začetne vrednosti \(z_0 \in \CC\). Takšno zaporedje imenujemo \emph{orbita} števila \(z_0\). Dokazali bomo naslednji izrek.

\begin{izrek}[Orbite kompleksne eksponentne preslikave] \label{thm:orbits}
    Vsaka od naslednjih množic je za preslikavo \(f \colon \CC \to \CC\); \(z \mapsto e^{z}\) gosta v kompleksni ravnini:
    \begin{enumerate}
        \item množica števil, katerih orbita divergira k \(\infty\);
        \item množica števil, katerih orbita gosto pokrije kompleksno ravnino;
        \item množica periodičnih točk, t.j.~števil \(z_0\), za katere obstaja \(k \in \NN\), da je \(z_k = z_0\).
    \end{enumerate}
\end{izrek}

\noindent Iz izreka nemudoma sledi, da je eksponentna preslikava \emph{kaotična}, kar bomo natančno definirali v naslednjem poglavju. V njem bomo obravnavali osnovno teorijo dinamičnih sistemov na metričnih prostorih. Nato v poglavju \ref{sec:hipgeom} ponovimo osnove kompleksne analize ter na kompleksnih ravnini uvedemo hiperbolično metriko, katere lastnosti nam bodo koristile pri dokazu osrednjega izreka. Z rezultati kompleksne analize v poglavju \ref{sec:komdin} dodatno razvijemo teorijo dinamičnih sistemov na kompleksni ravnini. Nato sledi dokaz izreka \ref{thm:orbits} v poglavju \ref{sec:dokaz}. Ta bo služil kot dober uvod v \emph{transcendentno dinamiko}, ki jo bomo dodatno obravnavali v zadnjem poglavju \ref{sec:druzina}.
